% !TEX root = ../my-thesis.tex

\chapter{Grundlagen und verwandte Arbeiten}
\label{sec:basics}



\section{Wissenschaftliche Workflows und Workflow-Management-Systeme}
\label{sec:basics:wms}
Wissenschaftliche Workflows beschreiben wissenschaftliche
Berechnungs- oder Analyseprozesse, die aus mehreren miteinander
verbundenen Rechenschritten bestehen. Ein solcher Workflow verfolgt
ein bestimmtes Forschungs- oder Analyseziel und beschreibt, wie die
einzelnen Schritte innerhalb des Ablaufs zusammenhängen. Dazu
gehören insbesondere ihre Reihenfolge sowie die zwischen ihnen
bestehenden Daten- und
Kontrollabhängigkeiten~\cite[S.~2]{suter2026terminology}.
Anschaulich lässt sich diese Struktur als Graph darstellen. Die
Knoten stehen dabei für die einzelnen Rechenschritte, in der
Literatur meist als Tasks bezeichnet, während die Kanten die Daten-
oder Kontrollabhängigkeiten zwischen ihnen
darstellen~\cite[S.~28\,f.]{badia2017workflows}. Eine verbreitete
Form eines Rechenschritts ist ein eigenständiges ausführbares
Programm, das Eingabedaten verarbeitet und daraus Ausgaben
erzeugt~\cite[S.~4]{suter2026terminology}. Die zwischen den
Rechenschritten ausgetauschten Daten liegen dabei häufig als Dateien
vor~\cite[S.~29]{badia2017workflows}. In dieser
Arbeit wird durchgehend der Begriff Schritt verwendet.
 
Die Ausführung eines Workflows wird durch ein
Workflow-Management-System gesteuert. Badia et al. definieren ein
solches System als Softwareumgebung, die die Ausführung einer Menge
voneinander abhängiger Rechenaufgaben orchestriert, welche
untereinander Daten austauschen, um ein bestimmtes Experiment
durchzuführen~\cite[S.~28]{badia2017workflows}.
Zu seinen Aufgaben gehören je nach System die Planung und
Koordination der Ausführung, die Übertragung von Daten sowie die
Überwachung des Ausführungszustands~\cite[S.~29]{badia2017workflows}\cite[S.~2]{suter2026terminology}.

\section{Workflow-Struktur und Ausführung}
\label{sec:basics:struktur}
Zwischen den Schritten eines Workflows bestehen Abhängigkeiten, die
bestimmen, wann ein Schritt ausgeführt werden kann. Badia et al.
unterscheiden dabei Datenabhängigkeiten von
Kontrollabhängigkeiten~\cite[S.~28]{badia2017workflows}. Eine
Datenabhängigkeit besteht, wenn ein Schritt Daten benötigt, die
von einem anderen Schritt erzeugt werden. Eine
Kontrollabhängigkeit besteht dagegen, wenn der Abschluss eines
Schritts Voraussetzung dafür ist, dass ein anderer Schritt
ausgeführt werden kann~\cite[S.~259]{zhou2008control}. Das
Workflow-Management-System orchestriert die Ausführung so, dass die
bestehenden Abhängigkeiten eingehalten
werden~\cite[S.~4]{suter2026terminology}.
 
Aus den Abhängigkeiten ergibt sich, welche Schritte nacheinander und
welche gleichzeitig ausgeführt werden können. Bei einer sequenziellen Ausführung werden aufeinanderfolgende Schritte einer nach dem anderen
ausgeführt~\cite[S.~8]{russell2006patterns}.
Ein Ausführungszweig kann sich aber auch in mehrere Zweige
aufteilen, die gleichzeitig ausgeführt
werden~\cite[S.~9]{russell2006patterns}. Mehrere Zweige können anschließend wieder in einem gemeinsamen Folgeschritt zusammengeführt werden, der erst beginnt, wenn alle eingehenden Zweige abgeschlossen sind~\cite[S.~10]{russell2006patterns}.

Wissenschaftliche Workflows können aus solchen einfachen Strukturen
zu komplexeren Abläufen zusammengesetzt
sein~\cite[S.~2]{bharathi2008characterization}. Bharathi et al.
beschreiben als einfachste Struktur einen einzelnen
Verarbeitungsschritt, der Eingabedaten verarbeitet und daraus eine
Ausgabe erzeugt~\cite[S.~3]{bharathi2008characterization}. Werden mehrere solcher Schritte hintereinandergeschaltet, entsteht eine Pipeline, in der jeder Schritt auf der Ausgabe des vorherigen
arbeitet und seine Ausgabe an den nächsten weitergibt~\cite[S.~3]{bharathi2008characterization}.
Daneben beschreiben sie Schritte, die einen Datensatz in kleinere
Teilstücke zerlegen, damit nachfolgende Schritte diese parallel
verarbeiten können, sowie Schritte, die die Ergebnisse mehrerer
Schritte zu einem gemeinsamen Ergebnis
zusammenführen~\cite[S.~3]{bharathi2008characterization}. Suter et
al. nennen Datenverarbeitungs-Pipelines und Fan-out/Fan-in-Muster
entsprechend als einfache Formen wissenschaftlicher
Workflows~\cite[S.~4]{suter2026terminology}. Die Verbindung aus
Zerlegen, paralleler Verarbeitung der Teilstücke und Zusammenführen
wird in dieser Arbeit als Scatter-Gather-Muster bezeichnet. Die
parallel ausgeführten Einheiten desselben Schritts werden im
Folgenden als Verarbeitungsinstanzen bezeichnet.



\section{HPC-Systeme und Ressourcenverwaltung}
\label{sec:basics:hpc}
High-Performance Computing (HPC) dient der Bearbeitung
rechenintensiver Probleme, deren Berechnung auf einem einzelnen
Rechner zu aufwendig oder zu zeitintensiv wäre. Dafür kommen unter anderem Supercomputer zum Einsatz, die aus hunderten oder tausenden über ein
Hochgeschwindigkeitsnetzwerk verbundenen Rechnern bestehen
können~\cite{zdvhpc}. Der Zugang zu einem HPC-System erfolgt über
Login-Knoten, auf denen Nutzende Daten bereitstellen,
Rechenaufträge vorbereiten und einreichen sowie deren Zustand
verfolgen können~\cite{slurmdocs-admin}. Die eigentlichen
Berechnungen werden dagegen auf den Rechenknoten
ausgeführt~\cite{slurmdocs-admin}. Da mehrere Aufträge um die
verfügbaren Rechenressourcen konkurrieren können, muss deren
Zuteilung koordiniert werden. Ein Ressourcenverwalter teilt dazu
den Zugriff auf Rechenressourcen zu und verwaltet konkurrierende
Anforderungen~\cite{slurmdocs-quickstart}.
Workflow-Management-Systeme können die Ressourcenzuteilung an einen
solchen Ressourcenverwalter delegieren. In HPC-Systemen übernimmt
diese Aufgabe typischerweise ein
Batch-Scheduler, der eingereichte
Rechenaufträge zur Ausführung
einplant~\cite[S.~5]{suter2026terminology}. In dieser Arbeit wird hierfür Slurm eingesetzt, ein System zur Ressourcenverwaltung und Job-Steuerung auf Linux-Clustern~\cite{slurmdocs-quickstart}.

Rechenaufträge werden bei Slurm als Jobs eingereicht. Für
Batch-Jobs erfolgt dies mit dem Kommando \texttt{sbatch}, das ein
Jobskript an Slurm übermittelt. Das Skript kann über
\texttt{\#SBATCH}-Direktiven unter anderem die benötigten CPUs, den
Hauptspeicher, ein Zeitlimit und die gewünschte Partition
festlegen~\cite{slurmdocs-sbatch}. Nach erfolgreicher Einreichung
weist Slurm dem Job eine Job-ID zu. Stehen die angeforderten
Ressourcen nicht unmittelbar zur Verfügung, kann der Job zunächst
in der Warteschlange ausstehender Jobs verbleiben, bis die
benötigten Ressourcen verfügbar
sind~\cite{slurmdocs-sbatch}. Die Rechenknoten sind in Slurm in
Partitionen gegliedert, also in Gruppen von Knoten, innerhalb derer
den Jobs Ressourcen zugeteilt
werden~\cite{slurmdocs-quickstart}. Teilt Slurm einem Job die
angeforderten Ressourcen zu, erhält er eine Allokation, also den
Zugriff auf diese Ressourcen für eine begrenzte
Dauer~\cite{slurmdocs-quickstart}. Der Zugriff kann exklusiv oder
geteilt mit anderen Jobs erfolgen. Bei einer exklusiven Allokation
teilt der Job die zugeteilten Knoten nicht mit anderen laufenden
Jobs~\cite{slurmdocs-sbatch}. Nach erteilter Allokation führt Slurm
das Jobskript auf dem ersten der zugeteilten Knoten
aus~\cite{slurmdocs-sbatch}.

Vor der eigentlichen Ausführung eines Jobs können auf den
zugeteilten Rechenknoten systemseitige Vorbereitungen stattfinden.
Slurm stellt dafür unter anderem Prolog-Programme bereit, die auf
den Rechenknoten ausgeführt werden, bevor dort die Aufgaben des
Jobs beginnen~\cite{slurmdocs-prolog}. Diese vorbereitenden
Arbeiten sind von der eigentlichen Nutzlast des Jobs zu
unterscheiden. Slurm überträgt ausführbare Dateien und Daten nicht
automatisch auf die einem Job zugeteilten Knoten. Die benötigten
Dateien müssen daher entweder auf dem lokalen Speicher des
jeweiligen Knotens oder in einem gemeinsam erreichbaren
Dateisystem vorliegen~\cite{slurmdocs-quickstart}. Werden
aufeinanderfolgende Schritte eines Workflows auf unterschiedlichen
Knoten ausgeführt, müssen die zwischen ihnen weitergegebenen
Dateien für die jeweils nachfolgenden Knoten erreichbar sein,
beispielsweise über ein gemeinsames Dateisystem.

\section{Untersuchte Workflow-Management-Systeme}
\label{sec:basics:systeme}


\subsection{Nextflow}
\label{sec:basics:nextflow}
Nextflow ist ein Workflow-Management-System aus der Bioinformatik.
Es wurde entwickelt, um rechnergestützte Analysen über
unterschiedliche Rechenumgebungen hinweg reproduzierbar zu machen,
und setzt dafür unter anderem auf
Container~\cite[S.~316]{ditommaso2017nextflow}. Workflows werden in
einer systemeigenen Sprache beschrieben, die auf dem
Datenfluss-Modell beruht. Ein Workflow setzt sich aus Prozessen
zusammen, die über Channels verbunden werden, indem die Ausgaben
eines Prozesses als Eingaben eines anderen dienen. Ein Prozess
startet, sobald seine Eingaben über die Channels
eintreffen~\cite[S.~316\,f.]{ditommaso2017nextflow}.


\subsection{StreamFlow}
\label{sec:basics:streamflow}
StreamFlow ist ein Workflow-Management-System, das die Ausführung
eines Workflows über mehrere, auch unterschiedliche Ausführungsumgebungen hinweg ermöglicht, etwa über Cloud- und HPC-Ressourcen gemeinsam~\cite[S.~1723]{colonnelli2021streamflow}. Dazu erweitert StreamFlow ein klassisches Workflow-Modell um die Beschreibung verschiedener Ausführungsumgebungen und die Zuordnung von Workflow-Schritten zu diesen Umgebungen~\cite[S.~1723]{colonnelli2021streamflow}. Als Workflow-Beschreibung verwendet StreamFlow die Common Workflow Language (CWL)~\cite[S.~1723]{colonnelli2021streamflow}.

\subsection{Merlin}
\label{sec:basics:merlin}
Merlin ist ein Workflow-Management-System aus dem HPC-Umfeld. Es wurde
entwickelt, um sehr große Ensembles von Simulationsläufen zu
erzeugen, etwa als Trainingsdaten für maschinelles
Lernen~\cite[S.~255\,f.]{peterson2022merlin}. Workflows werden in
einer Spezifikation beschrieben, in der eine Studie aus Schritten
mit Kommandos, Abhängigkeiten und Parametern festgelegt
wird~\cite[S.~258]{peterson2022merlin}.
Für die Ausführung verwendet Merlin Celery, ein verteiltes System zur
Aufgabenverwaltung~\cite[S.~258]{peterson2022merlin}\cite{merlindocs-configuration}. Merlin übersetzt die Schritte des Workflows dazu in Aufgaben und stellt diese über einen Nachrichtenvermittler (Broker) in einer
Warteschlange bereit~\cite[S.~258]{peterson2022merlin}\cite{merlindocs-configuration}. Worker sind Prozesse, die Aufgaben aus dieser Warteschlange
beziehen und ausführen~\cite[S.~258]{peterson2022merlin}\cite{merlindocs-configuration}. Dadurch sind die Worker von der Erzeugung der Aufgaben entkoppelt und können unabhängig davon auf den Rechenressourcen gestartet
werden~\cite[S.~261]{peterson2022merlin}.


\subsection{Pegasus}
\label{sec:basics:pegasus}
Pegasus beruht auf der Trennung zwischen der Beschreibung eines
Workflows und seiner konkreten Ausführung. Nutzende erstellen
einen abstrakten Workflow, den Pegasus in einen ausführbaren
Workflow überführt~\cite[S.~18]{deelman2015pegasus}. Der abstrakte
Workflow wird als gerichteter azyklischer Graph in einem
XML-Format namens DAX beschrieben~\cite[S.~19]{deelman2015pegasus}. Die Überführung in einen ausführbaren Workflow übernimmt ein Planer, der Mapper.
Dabei bestimmt er unter anderem, auf welchen Ressourcen die Jobs
ausgeführt und von welchen Speicherorten benötigte Daten bezogen
werden~\cite[S.~20]{deelman2015pegasus}. Für die anschließende
Ausführung verwendet Pegasus standardmäßig HTCondor DAGMan als
Workflow-Engine. DAGMan berücksichtigt die Abhängigkeiten des
ausführbaren Workflows und gibt einen Job zur Ausführung frei,
sobald dessen Vorgänger abgeschlossen
sind~\cite[S.~20]{deelman2015pegasus}. Die einzelnen Jobs werden
dabei über HTCondor verwaltet und können je nach
Ausführungsumgebung lokal oder an entfernte Ressourcenverwalter
übergeben werden~\cite[S.~20]{deelman2015pegasus}.

\subsection{AiiDA}
\label{sec:basics:aiida}
AiiDA stellt die Nachvollziehbarkeit der erzeugten Daten in den
Mittelpunkt. Provenance bezeichnet dabei die Herkunft eines
Datenelements und den Vorgang, durch den es entstanden
ist~\cite[S.~316]{buneman2001provenance}. Jede vom System
ausgeführte Berechnung und jeder Workflow werden automatisch in
einem Provenance-Graphen aufgezeichnet, in dem Prozesse und Daten
als Knoten und ihre Beziehungen als Kanten festgehalten werden, um
die Reproduzierbarkeit der Ergebnisse zu
ermöglichen~\cite[S.~2]{huber2020aiida}. Für die Ablage stellt
AiiDA zwei Speicher bereit: Die Eigenschaften der Knoten liegen
als Attribute in einer relationalen Datenbank, Dateien werden in
einem Dateirepository abgelegt~\cite[S.~8]{huber2020aiida}.
Workflows werden in Python beschrieben. Dazu stellt AiiDA die
Klasse \texttt{WorkChain} bereit, in der die Schritte als
logischer Ablauf in einer Prozessspezifikation festgelegt
werden~\cite[S.~4\,f.]{huber2020aiida}. Berechnungen, die externe
Programme ausführen, etwa über einen Job-Scheduler, übernimmt die
Klasse \texttt{CalcJob}~\cite[S.~4]{huber2020aiida}.
Für die Ausführung verwendet AiiDA mehrere Runner, die die
eingereichten Prozesse abarbeiten. Sie werden von einem Daemon
überwacht, einem dauerhaft im Hintergrund laufenden
Prozess~\cite[S.~3]{huber2020aiida}. Die dabei
anfallenden Aufgaben werden über eine Aufgabenwarteschlange
verteilt, die mit dem Nachrichtenvermittler RabbitMQ umgesetzt
ist~\cite[S.~3]{huber2020aiida}.

\section{Verwandte Arbeiten und Abgrenzung}
\label{sec:basics:related}
Workflow-Management-Systeme wurden bereits in verschiedenen
Arbeiten anhand unterschiedlicher Kriterien verglichen.
Diercks et al. leiten Anforderungen an Workflow-Management-Systeme
aus Anwendungsszenarien des Computational Science and Engineering
ab und vergleichen sechs Systeme anhand dieser Anforderungen.
Im Mittelpunkt stehen dabei insbesondere die Beschreibung,
Reproduzierbarkeit und Wiederverwendung wissenschaftlicher
Workflows~\cite[S.~1]{diercks2023evaluation}.
Singh et al. vergleichen fünf wissenschaftliche
Workflow-Management-Systeme anhand qualitativer Merkmale wie
Orchestrierung, Datenverwaltung und Provenance.
Dabei betrachten sie die Systeme insbesondere im Hinblick auf
daten- und rechenintensive wissenschaftliche
Workflows~\cite[S.~4553]{singh2019evaluating}.
Larsonneur et al. vergleichen fünf Workflow-Management-Systeme,
darunter Nextflow und Pegasus, anhand eines
Bioinformatik-Workflows. Die Systeme werden dabei anhand von
zehn Effizienzmetriken wie Laufzeit, CPU- und Speicherbedarf
bewertet~\cite[S.~2773\,f.]{larsonneur2018evaluating}.

Die vorliegende Arbeit unterscheidet sich von diesen Arbeiten
durch die Auswahl der Systeme und den experimentellen Aufbau.
Mit StreamFlow und Merlin werden zwei Systeme betrachtet, die in
den genannten Vergleichen nicht vertreten sind. Für alle fünf
Systeme werden dieselben Workflow-Muster mit identischer Nutzlast
sowohl lokal als auch auf einem HPC-System mit Slurm betrachtet.